% ============================================================================
%  A/01-nghia-nam-o-dau — file chương
%  Ngày tạo: 2026-09-06 | Sửa gần nhất: 2026-09-06 | Ôn gần nhất: chưa
%  Dependency: không. Mức độ: cốt lõi.
% ============================================================================
\documentclass[../../english.tex]{subfiles}
\begin{document}
\chapter{Nghĩa nằm ở đâu (Where Meaning Lives)}
\ptrtarget{A/01}\ptrtarget{A/01-nghia-nam-o-dau}
\dependency{không --- đây là chương mở đầu. Mọi chương sau đều dựa trên cách hiểu về nghĩa được dựng ở đây.}

\begin{tomtat}
Chương này thay một giả định ngầm mà gần như mọi người học tiếng Anh đều mang theo: rằng mỗi từ có một nghĩa, và tra từ điển là lấy được nghĩa đó. Thực tế thì từ điển chỉ ghi lại \emph{cách dùng phổ biến nhất} của một từ trong kho văn bản thật, và nghĩa cụ thể trong một câu cụ thể được quyết định bởi ba tầng: nghĩa quy chiếu, sắc thái, và ý định của người viết. Chương giới thiệu ba tầng đó, chỉ ra vì sao hai từ ``đồng nghĩa'' hầu như không bao giờ thay thế được cho nhau, và cho bốn phép kiểm tra để biết bạn đã hiểu đúng một câu hay chưa. Nếu chỉ nhớ một điều: bạn biết một từ khi bạn biết nó \emph{thường đi với những từ nào} --- và đó là thứ từ điển song ngữ không cho bạn.
\end{tomtat}

\begin{nentang}
\kn{nghĩa quy chiếu (denotation)}{phần nghĩa ``khách quan'': từ đó chỉ vào sự vật, khái niệm hay quan hệ nào. \emph{Cheap} và \emph{inexpensive} có cùng nghĩa quy chiếu: giá thấp.}
\kn{sắc thái (connotation)}{phần nghĩa mang thái độ, đánh giá, hoặc gợi liên tưởng. \emph{Cheap} gợi thêm ``rẻ tiền, kém chất lượng''; \emph{inexpensive} thì không.}
\kn{văn vực (register)}{mức trang trọng và bối cảnh dùng của một từ: học thuật, trung tính, đời thường, kỹ thuật. Dùng sai văn vực làm câu ``không sai'' mà vẫn lạc lõng.}
\kn{đa nghĩa (polysemy)}{một từ có nhiều nghĩa liên quan nhau; ngữ cảnh chọn ra nghĩa đúng. \emph{Set} trong \emph{a set of points} và trong \emph{set a threshold} là hai nghĩa khác nhau.}
\kn{kết hợp từ (collocation)}{những từ thường xuyên đi cùng nhau: \emph{make a decision} chứ không \emph{do a decision}. Chi tiết ở \ptr{D/17}.}
\kn{ngữ dụng (pragmatics)}{nghiên cứu về điều người nói \emph{ngụ ý} nhưng không nói thẳng. Chi tiết ở \ptr{B/08}.}
\end{nentang}

\begin{khoigoi}
\h{Bạn hiểu mọi từ trong một câu nhưng vẫn không chắc tác giả khẳng định điều gì. Cái thiếu là gì, nếu không phải từ vựng?}
\h{Vì sao \emph{cheap} và \emph{inexpensive} không thay thế được cho nhau, dù từ điển ghi cùng một nghĩa tiếng Việt?}
\h{Từ điển được biên soạn thế nào, và điều đó nói gì về cách bạn nên đọc một mục từ?}
\h{Trong bài báo khoa học, cùng một từ như \emph{significant} hay \emph{robust} lại có nghĩa khác nhau giữa các ngành. Phải đọc thế nào để không hiểu nhầm?}
\h{Làm sao \emph{kiểm tra} rằng mình đã hiểu đúng một câu khó, thay vì chỉ cảm thấy đã hiểu?}
\end{khoigoi}

Hãy bắt đầu bằng một cảnh rất cụ thể. Bạn đọc trong phần thảo luận của một bài báo:

\begin{quote}
\emph{These results suggest that the proposed method may offer a modest improvement over existing approaches, at least in the settings we considered.}
\end{quote}

Mọi từ đều quen. Bạn dịch được từng chữ. Nhưng câu này thực ra đang nói gì? Tác giả \emph{có} khẳng định phương pháp của họ tốt hơn không? Mạnh tới mức nào? Và vì sao câu lại dài đến vậy trong khi ý chính chỉ là ``phương pháp tốt hơn một chút''?

Câu trả lời là: gần như mọi từ thừa trong câu đó đều mang thông tin. \emph{Suggest} chứ không phải \emph{show} hay \emph{prove} --- tác giả nói rằng dữ liệu \emph{gợi ý}, không phải chứng minh. \emph{May} thêm một lớp dè dặt nữa. \emph{Modest} thừa nhận cải thiện nhỏ. \emph{At least in the settings we considered} là một hàng rào rõ ràng: đừng suy rộng ra ngoài các tình huống chúng tôi đã thử. Người đọc quen với văn khoa học tiếng Anh đọc câu này ra một thông điệp rất khác so với ``phương pháp của chúng tôi tốt hơn'': họ đọc ra một khẳng định \emph{đã được hiệu chỉnh cẩn thận}, và họ biết rằng nếu tác giả có bằng chứng mạnh hơn thì đã viết mạnh hơn.

Đó là điều mà từ điển không dạy được, và đó là nội dung của cả quyển sách này.

\section{Từ điển cho bạn cái gì, và không cho cái gì}

Có một hiểu lầm phổ biến về từ điển: rằng nó ghi lại \emph{nghĩa đúng} của từ, như một quy định. Thực tế thì các từ điển hiện đại làm việc theo hướng ngược lại: người biên soạn quan sát hàng trăm triệu câu tiếng Anh thật trong một kho ngữ liệu, xem mỗi từ thực sự được dùng thế nào, rồi \emph{mô tả} lại cách dùng phổ biến nhất. Từ điển là bản tóm tắt thống kê của thói quen sử dụng, không phải bộ luật.

Hệ quả rất thực dụng. Thứ nhất, thứ tự các nghĩa trong một mục từ thường phản ánh tần suất, nên nghĩa đầu tiên là nghĩa hay gặp nhất chứ không phải nghĩa ``gốc''. Thứ hai, phần ví dụ trong mục từ quan trọng \emph{hơn} phần định nghĩa, vì nó cho bạn thấy từ đó đi với những từ nào và ở văn vực nào. Thứ ba, một mục từ tốt luôn ghi thêm những thứ mà từ điển song ngữ bỏ qua: giới từ đi kèm, danh từ đếm được hay không, và các cụm cố định.

Từ đó rút ra một thay đổi thói quen đáng giá: \emph{dùng từ điển đơn ngữ dành cho người học} thay vì từ điển Anh--Việt, ít nhất cho những từ bạn định dùng lại. Từ điển học tập định nghĩa bằng ngôn ngữ đơn giản và luôn kèm ví dụ thật; từ điển song ngữ cho bạn một nhãn tiếng Việt và lấy đi toàn bộ phần còn lại. Hai công cụ này phục vụ hai mục đích khác nhau, và nhầm lẫn giữa chúng là lý do rất nhiều người có vốn từ lớn nhưng viết ra vẫn không tự nhiên.

\begin{cambay}
Bẫy nguy hiểm nhất với người học: dùng từ điển đồng nghĩa để ``làm câu văn phong phú hơn''. Từ điển đồng nghĩa liệt kê những từ có nghĩa quy chiếu gần nhau, nhưng nó không nói cho bạn sắc thái, văn vực, hay kết hợp từ. Thay \emph{use} bằng \emph{utilise} không làm câu học thuật hơn --- nó chỉ làm câu nặng hơn và thường sai văn vực. Quy tắc an toàn: chỉ dùng một từ đồng nghĩa sau khi đã \emph{nhìn thấy} nó trong ít nhất ba câu thật thuộc đúng loại văn bản bạn đang viết.
\end{cambay}

\section{Ba tầng nghĩa}

Khi đọc một câu, bạn thực ra đang phục dựng ba tầng chồng lên nhau. Tách chúng ra giúp bạn biết mình đang thiếu tầng nào.

\emph{Tầng một --- nghĩa quy chiếu.} Từ này chỉ vào cái gì? Đây là tầng mà từ điển phục vụ tốt, và cũng là tầng duy nhất mà phần lớn người học chú ý.

\emph{Tầng hai --- sắc thái và văn vực.} Từ này mang thái độ gì, và thuộc loại văn bản nào? Đây là chỗ \emph{cheap} khác \emph{inexpensive}, \emph{stubborn} khác \emph{persistent}, \emph{cause} khác \emph{lead to}. Trong văn khoa học, tầng này thường quyết định mức độ mạnh yếu của một khẳng định.

\emph{Tầng ba --- ý định.} Người viết dùng câu này để \emph{làm gì}: khẳng định, dè dặt, phản bác nhẹ nhàng, nhường bộ trước khi phản đối, hay báo hiệu chuyển ý? Tầng này là ngữ dụng, và nó được nói kỹ ở \ptr{B/08}.

Bảng dưới lấy một nhóm động từ mà người viết bài báo dùng hằng ngày. Chúng có nghĩa quy chiếu gần nhau --- ``đưa ra một mệnh đề'' --- nhưng khác nhau hoàn toàn ở mức độ cam kết, và chọn sai từ làm người đọc hiểu sai độ mạnh của kết quả bạn báo cáo.

\begin{tuvung}[Cùng chỉ ``nói ra một điều'', khác nhau ở mức cam kết]
\tv{show}{cho thấy, có bằng chứng trực tiếp}{Mạnh; ngụ ý dữ liệu trong bài đủ để kết luận}{Figure 3 \textbf{shows} that the error decreases with more samples.}
\tv{demonstrate}{chứng tỏ, thường qua thí nghiệm có kiểm soát}{Mạnh hơn \emph{show} một chút, trang trọng hơn}{We \textbf{demonstrate} the approach on three public datasets.}
\tv{suggest}{gợi ý, chưa kết luận}{Dè dặt; dùng khi bằng chứng có nhưng chưa dứt khoát}{These results \textbf{suggest} that latency, not bandwidth, is the bottleneck.}
\tv{indicate}{chỉ dấu, thiên về dữ liệu}{Trung tính, hơi mạnh hơn \emph{suggest}}{The logs \textbf{indicate} a memory leak in the tracking thread.}
\tv{argue}{lập luận, đưa quan điểm}{Ngụ ý đây là \emph{quan điểm} cần bảo vệ, không phải dữ kiện}{We \textbf{argue} that the standard benchmark is misleading.}
\tv{claim}{tuyên bố}{Thường mang chút hoài nghi khi nói về người khác}{Prior work \textbf{claims} real-time performance but reports no timings.}
\tv{prove}{chứng minh}{Rất mạnh; trong ngữ cảnh toán học có nghĩa hẹp và chặt}{We \textbf{prove} that the algorithm terminates in \(O(n\log n)\) steps.}
\end{tuvung}

Điều đáng chú ý ở bảng trên: \emph{claim} khi nói về công trình của người khác gần như luôn mang sắc thái hoài nghi nhẹ --- người đọc hiểu ngầm rằng tác giả không hoàn toàn tin lời tuyên bố đó. Đây chính là tầng ba: cùng một sự việc, chọn \emph{claim} thay vì \emph{show} là một hành động, không chỉ là một lựa chọn từ.

\section{Ngữ cảnh quyết định nghĩa}

Có một câu nói của nhà ngôn ngữ học J.\,R.\ Firth năm 1957 đã trở thành nền tảng của cả ngành từ vựng học hiện đại: bạn sẽ biết một từ qua những từ đi cùng với nó\cite{firth1957}. Nghe như một câu khẩu hiệu, nhưng nó có ba hệ quả rất cụ thể.

\emph{Hệ quả thứ nhất: ngữ cảnh chọn nghĩa.} Từ \emph{run} trong từ điển có hàng chục nghĩa, nhưng bạn không bao giờ phải cân nhắc chúng khi đọc \emph{run an experiment} hay \emph{the algorithm runs in linear time}. Các từ xung quanh đã loại bỏ mọi nghĩa khác trước khi bạn kịp nghĩ. Vì vậy, học một từ trong câu luôn hiệu quả hơn học nó trong danh sách --- không phải vì dễ nhớ hơn, mà vì bạn học được cả cơ chế chọn nghĩa.

\emph{Hệ quả thứ hai: có những từ mang sẵn một ``bầu không khí''.} Các nhà nghiên cứu ngữ liệu gọi hiện tượng này là sắc thái ngữ nghĩa\cite{louw1993}: một số từ hầu như chỉ xuất hiện trong ngữ cảnh tiêu cực dù bản thân chúng trung tính. \emph{Cause} là ví dụ điển hình --- trong kho ngữ liệu, nó đi với \emph{damage, problems, delay, failure} nhiều hơn hẳn với những thứ tốt đẹp; khi muốn nói về kết quả tích cực, người bản ngữ thường viết \emph{lead to} hoặc \emph{result in}. Người học không biết điều này viết \emph{This caused a significant improvement} --- ngữ pháp đúng, nhưng câu nghe lệch.

\emph{Hệ quả thứ ba: nghĩa của một từ có thể trôi theo ngành.} Trong thống kê, \emph{significant} có nghĩa kỹ thuật hẹp về mức ý nghĩa; trong đời thường nó chỉ là ``đáng kể''. Trong học máy, \emph{robust} thường có nghĩa ``chịu được nhiễu và dữ liệu ngoại lai''; trong thống kê cổ điển nó có định nghĩa hẹp hơn; trong kỹ thuật phần mềm nó lại gần với ``không sập''. Đây là câu trả lời cho câu hỏi mở đầu thứ tư, và quy tắc thực hành rất đơn giản: \emph{trong tài liệu chuyên ngành, hãy tìm định nghĩa mà chính tác giả đưa ra trước khi dùng nghĩa thường ngày}. Bài báo tốt luôn định nghĩa các thuật ngữ then chốt ở phần đầu; nếu không, đó cũng là một thông tin về chất lượng bài.

\begin{trucgiac}
Một cách nghĩ hữu ích cho người làm kỹ thuật: từ trong tiếng Anh giống \emph{tên hàm bị nạp chồng} trong lập trình. Cùng một tên, nhiều phiên bản, và trình biên dịch chọn phiên bản đúng dựa trên kiểu của các đối số. Ở đây ``kiểu của đối số'' chính là các từ xung quanh. Điều đó giải thích vì sao học từ rời rạc lại kém hiệu quả: bạn đang nhớ tên hàm mà không nhớ chữ ký của nó, nên khi cần gọi thì không biết truyền gì vào.
\end{trucgiac}

\section{Bốn phép kiểm tra đã hiểu đúng}

Cảm giác ``tôi hiểu rồi'' rất dễ đánh lừa, vì nó chỉ đòi hỏi các từ quen thuộc chứ không đòi hỏi bạn tái tạo được nội dung. Bốn phép kiểm tra dưới đây rẻ và bắt được phần lớn hiểu nhầm.

\emph{Phép một --- diễn đạt lại.} Nói lại câu đó bằng từ của bạn, không nhìn bản gốc, và không dùng lại các từ khoá. Nếu bạn chỉ đảo được trật tự từ mà không thay được từ, bạn chưa hiểu, chỉ mới nhận ra.

\emph{Phép hai --- thay thế.} Thử thay từ then chốt bằng một từ ``đồng nghĩa'' và hỏi câu có đổi nghĩa không. Với ví dụ ở đầu chương: thay \emph{suggest} bằng \emph{show} --- câu mạnh hẳn lên. Vậy \emph{suggest} đang gánh một phần nghĩa quan trọng, và bạn vừa xác định được nó gánh phần nào.

\emph{Phép ba --- phủ định.} Hỏi ``nếu điều này sai thì thế giới sẽ khác thế nào?''. Một câu mà bạn không hình dung được phiên bản phủ định của nó thì thường là câu bạn chưa hiểu --- hoặc là câu rỗng, điều cũng đáng biết.

\emph{Phép bốn --- xác định phạm vi.} Câu này áp dụng cho cái gì và \emph{không} áp dụng cho cái gì? Trong văn khoa học, phạm vi thường nằm ở các cụm bị bỏ qua khi đọc nhanh: \emph{in this setting}, \emph{for small \(n\)}, \emph{under the assumption that\ldots} Bỏ qua chúng là cách phổ biến nhất để hiểu một bài báo mạnh hơn ý tác giả.

\begin{cannho}
Ba câu hỏi nên thành phản xạ khi đọc một câu quan trọng: (i) tác giả cam kết \emph{mạnh tới đâu} --- \emph{show}, \emph{suggest}, hay \emph{may}; (ii) phạm vi áp dụng là gì --- có cụm giới hạn nào không; (iii) nếu thay từ then chốt bằng từ gần nghĩa thì câu đổi thế nào. Ba câu này lấy ra phần lớn thông tin mà đọc lướt bỏ sót.
\end{cannho}

\section{Gốc rễ: từ tam giác nghĩa tới kho ngữ liệu}

Câu hỏi ``nghĩa là gì'' cổ như triết học, nhưng có ba mốc đã thay đổi hẳn cách chúng ta \emph{học} ngoại ngữ.

Mốc thứ nhất là ý tưởng rằng từ không nối thẳng vào sự vật. Năm 1923, Ogden và Richards đưa ra sơ đồ tam giác nghĩa: giữa \emph{ký hiệu} (từ) và \emph{sự vật} luôn có một đỉnh thứ ba là \emph{khái niệm trong đầu người dùng}\cite{ogden1923}. Hệ quả thực hành: hai người có thể dùng cùng một từ mà nghĩ về hai khái niệm khác nhau, và phần lớn hiểu nhầm trong đọc tài liệu kỹ thuật xảy ra ở đúng đỉnh này --- không phải vì bạn không biết từ, mà vì khái niệm bạn gán cho nó khác khái niệm của tác giả.

Mốc thứ hai là Firth với luận điểm rằng nghĩa được quyết định bởi các từ đi cùng\cite{firth1957}. Đây là bước chuyển từ việc coi nghĩa như một thuộc tính của từ sang coi nó như một hiện tượng \emph{thống kê} về cách từ được dùng.

Mốc thứ ba là khi luận điểm đó trở thành công cụ. Cuối thập niên 1980, dưới sự dẫn dắt của John Sinclair, một nhóm ở Đại học Birmingham dùng máy tính để phân tích một kho văn bản tiếng Anh thật rất lớn và biên soạn từ điển \emph{dựa trên} dữ liệu đó thay vì dựa trên trực giác của người biên soạn\cite{cobuild1987,sinclair1991}. Kết quả gây bất ngờ ngay cả với giới chuyên môn: nhiều nghĩa mà sách giáo khoa dạy là chính lại rất hiếm trong thực tế, còn nhiều cụm cố định mà không sách nào dạy thì xuất hiện liên tục. Chính từ đó ra đời loại từ điển dành cho người học có kèm ví dụ thật và thông tin về kết hợp từ --- công cụ mà mục 1 khuyên bạn dùng.

Bài học rút ra cho người học hôm nay khá trực tiếp: nếu nghĩa là hiện tượng thống kê về cách dùng, thì cách học hiệu quả nhất là \emph{gặp từ nhiều lần trong ngữ cảnh thật}, chứ không phải ghi nhớ định nghĩa. Toàn bộ phần E của quyển sách này được thiết kế quanh nhận xét đó.

\begin{gocre}
\gr{1923 --- Ogden \& Richards}{Vì sao hai người dùng cùng một từ mà hiểu khác nhau?}{Tam giác nghĩa: từ \(\rightarrow\) khái niệm \(\rightarrow\) sự vật. Hiểu nhầm thường nằm ở đỉnh khái niệm, không ở từ.}
\gr{1957 --- Firth}{Định nghĩa từ điển không giải thích được vì sao một số tổ hợp từ nghe tự nhiên còn một số thì không.}{``Biết một từ qua những từ đi cùng nó''; nghĩa trở thành hiện tượng thống kê về cách dùng.}
\gr{Cuối thập niên 1980 --- Sinclair và nhóm COBUILD}{Từ điển được viết theo trực giác người biên soạn, không theo dữ liệu thật.}{Từ điển biên soạn từ kho ngữ liệu; phát hiện rằng cụm cố định phổ biến hơn nhiều so với những gì sách giáo khoa dạy.}
\gr{1993 --- Louw và các nghiên cứu ngữ liệu}{Vì sao một số câu ngữ pháp đúng vẫn nghe lệch?}{Sắc thái ngữ nghĩa: nhiều từ trung tính lại hầu như chỉ xuất hiện trong ngữ cảnh tiêu cực (\emph{cause}, \emph{utterly}).}
\gr{Thập niên 2000--nay --- kho ngữ liệu mở}{Người học không tiếp cận được dữ liệu thật.}{Các kho tra cứu trực tuyến\cite{coca} cho phép bất kỳ ai kiểm tra ``người bản ngữ có thật sự viết thế này không'' trong vài giây (\ptr{E/20}).}
\end{gocre}

\section{Liên hệ ngang}

\begin{lienhe}
\lh{Tiếng Việt}{Từ đa nghĩa: ``chạy'' trong \emph{chạy bộ}, \emph{máy chạy}, \emph{chạy tiền}, \emph{chạy chương trình}}{Cùng cơ chế: ngữ cảnh chọn nghĩa trước khi bạn kịp cân nhắc. Điểm khác đáng chú ý: người Việt làm điều này tự động trong tiếng mẹ đẻ nhưng lại quên rằng tiếng Anh cũng vậy, nên hay ép một nghĩa duy nhất cho mỗi từ.}
\lh{Lập trình}{Nạp chồng hàm, phạm vi tên, không gian tên}{Không chỉ giống mà là mô hình rất sát: cùng một tên có nhiều nghĩa, và \emph{kiểu của các đối số} --- tức các từ xung quanh --- quyết định phiên bản nào được chọn. Học từ rời ngữ cảnh giống nhớ tên hàm mà quên chữ ký.}
\lh{Toán học và ký hiệu khoa học}{Một ký hiệu mang nghĩa khác nhau tuỳ ngành}{Cùng vấn đề với \emph{significant} hay \emph{robust}: nghĩa được \emph{quy ước} trong phạm vi một cộng đồng. Cách xử lý cũng giống nhau: tìm định nghĩa tác giả đưa ra trước khi áp nghĩa quen thuộc.}
\lh{Học máy}{Biểu diễn từ theo ngữ cảnh}{Các mô hình ngôn ngữ hiện đại học nghĩa của từ đúng bằng cách quan sát các từ xung quanh nó --- một hiện thực hoá kỹ thuật của luận điểm Firth năm 1957. Điều đó gợi ý ngược cho người học: dữ liệu ngữ cảnh chính là thứ mang thông tin về nghĩa.}
\lh{Nghiên cứu khoa học}{Định nghĩa thao tác trong một bài báo}{Khi tác giả viết ``in this paper, by \emph{accuracy} we mean\ldots'' thì họ đang làm đúng việc mà tam giác nghĩa gợi ý: cố định đỉnh khái niệm để tránh hiểu nhầm. Bạn nên làm điều tương tự khi viết (\ptr{C/14}).}
\lh{Luật và hợp đồng}{Điều khoản định nghĩa ở đầu văn bản}{Cùng cơ chế, mức độ nghiêm ngặt cao hơn: mọi thuật ngữ then chốt được định nghĩa trước, vì hậu quả của hiểu nhầm là pháp lý. Đọc hợp đồng tiếng Anh dạy rất nhanh thói quen đọc phạm vi và điều kiện.}
\end{lienhe}

\begin{traloi}
\tl{Cái thiếu không phải nghĩa quy chiếu của từng từ mà là hai tầng còn lại: sắc thái (từ này mạnh hay dè dặt, tích cực hay tiêu cực) và ý định (tác giả dùng câu này để làm gì --- khẳng định, hạn chế phạm vi, hay phản bác nhẹ). Trong ví dụ đầu chương, toàn bộ nội dung thật nằm ở \emph{suggest}, \emph{may}, \emph{modest} và cụm giới hạn cuối câu --- những chỗ mà người đọc chỉ chú ý nghĩa quy chiếu sẽ lướt qua.}
\tl{Vì chúng chỉ trùng nhau ở tầng nghĩa quy chiếu. \emph{Cheap} mang thêm sắc thái tiêu cực về chất lượng, nên \emph{a cheap solution} thường bị hiểu là giải pháp rẻ tiền theo nghĩa xấu, còn \emph{an inexpensive solution} chỉ nói về giá. Ngoài ra hai từ có kết hợp và văn vực khác nhau: \emph{inexpensive} trang trọng hơn và hầu như không dùng cho người. Không có hai từ nào trùng cả ba tầng --- nếu có, ngôn ngữ đã loại bỏ một trong hai.}
\tl{Từ điển hiện đại được biên soạn từ kho ngữ liệu: người biên soạn quan sát cách từ thực sự được dùng trong hàng trăm triệu câu rồi mô tả lại thói quen phổ biến nhất. Ba hệ quả cho cách đọc mục từ: thứ tự nghĩa phản ánh tần suất chứ không phải ``nghĩa gốc''; phần \emph{ví dụ} quan trọng hơn phần định nghĩa vì nó cho thấy kết hợp và văn vực; và những ghi chú về giới từ, tính đếm được, cụm cố định mới là phần bạn cần nhất khi định \emph{dùng} từ đó.}
\tl{Bằng cách đọc định nghĩa của chính tác giả trước khi áp nghĩa thường ngày. Trong văn khoa học, nhiều từ mang nghĩa được quy ước hẹp trong phạm vi một cộng đồng --- \emph{significant} trong thống kê, \emph{robust} trong học máy so với trong kỹ thuật phần mềm. Bài báo viết tốt sẽ định nghĩa các thuật ngữ then chốt ở phần đầu hoặc trích dẫn nguồn định nghĩa; nếu không có, hãy tìm cách họ \emph{đo} khái niệm đó --- định nghĩa thao tác thường nằm ở phần phương pháp.}
\tl{Bằng bốn phép kiểm ở mục 4: diễn đạt lại bằng từ khác (không chỉ đảo trật tự); thay từ then chốt bằng từ gần nghĩa và xem câu đổi thế nào; hình dung phiên bản phủ định của câu; và xác định rõ phạm vi --- câu này \emph{không} áp dụng cho cái gì. Phép thứ tư quan trọng đặc biệt khi đọc bài báo, vì các cụm giới hạn phạm vi là thứ hay bị bỏ qua nhất và là nguyên nhân chính khiến người đọc hiểu một kết quả mạnh hơn ý tác giả.}
\end{traloi}

\begin{dongian}
Nhiều người học tiếng Anh mang theo một hình dung sai: mỗi từ tiếng Anh là một cái nhãn dán lên một nghĩa tiếng Việt, và học tiếng Anh là học thuộc các cặp nhãn đó. Hình dung này giải thích vì sao đọc mãi vẫn không sâu.

Thực tế gần với chuyện khác. Một từ giống một người: bạn không hiểu người đó qua một câu mô tả, bạn hiểu qua việc thấy họ ở nhiều nơi, với những ai, trong hoàn cảnh nào. Từ cũng vậy --- \emph{cause} hầu như luôn xuất hiện cạnh những chuyện không hay, nên viết ``điều này gây ra một cải thiện lớn'' thì ngữ pháp đúng mà vẫn nghe kỳ, y như việc dùng một từ trang trọng trong câu nói đùa.

Có ba tầng bạn cần lấy ra từ mỗi câu quan trọng. Tầng thứ nhất: từ này chỉ cái gì. Tầng thứ hai: nó mang thái độ gì, trang trọng hay đời thường. Tầng thứ ba, và là tầng hay bị bỏ qua nhất: người viết dùng câu này để làm gì --- để khẳng định chắc chắn, hay để nói dè dặt, hay để rào trước một chỗ yếu.

Vì thế trong bài báo khoa học, khác nhau giữa ``chúng tôi cho thấy'' và ``kết quả gợi ý'' không phải chuyện văn phong. Đó là hai mức cam kết khác nhau, và người đọc có kinh nghiệm đọc ra ngay. Nếu một tác giả có bằng chứng mạnh, họ đã viết mạnh; khi họ viết dè dặt, đó là thông tin.

Cuối cùng, cách kiểm tra xem mình đã thật sự hiểu một câu: hãy thử nói lại nó bằng từ khác, và thử hỏi ``nếu câu này sai thì sao''. Nếu không làm được cả hai, cảm giác hiểu chỉ là cảm giác quen thuộc.
\motcau{Nghĩa không nằm trong từ điển; nó nằm ở chỗ từ đó thường xuất hiện.}
\chuadongian{Ranh giới giữa ``sắc thái'' và ``nghĩa'' không rạch ròi; với nhiều từ, chỉ có cách gặp nhiều lần trong văn bản thật mới cảm được.}
\end{dongian}

\begin{onbai}
\ar[3]{Ba tầng nghĩa}{Nghĩa quy chiếu (chỉ cái gì); sắc thái và văn vực (thái độ, mức trang trọng); ý định (người viết dùng câu này để làm gì).}
\ar[3]{Luận điểm Firth}{Biết một từ qua những từ đi cùng nó; nghĩa là hiện tượng thống kê về cách dùng, không phải thuộc tính cố định.}
\ar[3]{Từ điển thực ra là gì}{Bản mô tả thống kê thói quen sử dụng, không phải bộ luật; thứ tự nghĩa theo tần suất, và phần ví dụ quan trọng hơn phần định nghĩa.}
\ar[3]{Thang cam kết trong văn khoa học}{\emph{prove} > \emph{demonstrate} > \emph{show} > \emph{indicate} > \emph{suggest}; thêm \emph{may/might} để giảm mạnh. \emph{claim} về người khác mang hoài nghi nhẹ.}
\ar[3]{Bốn phép kiểm đã hiểu đúng}{Diễn đạt lại bằng từ khác; thay từ then chốt; hình dung phiên bản phủ định; xác định phạm vi áp dụng.}
\ar[3]{Bẫy từ điển đồng nghĩa}{Nó chỉ khớp nghĩa quy chiếu, bỏ mất sắc thái, văn vực và kết hợp từ. Chỉ dùng một từ sau khi đã thấy nó trong ít nhất ba câu thật cùng loại văn bản.}
\ar[2]{Sắc thái ngữ nghĩa}{Nhiều từ trung tính lại hầu như chỉ đi với ngữ cảnh tiêu cực: \emph{cause} đi với \emph{damage, delay, failure}; muốn nói kết quả tốt thì dùng \emph{lead to} hoặc \emph{result in}.}
\ar[2]{Từ chuyên ngành đổi nghĩa theo ngành}{\emph{significant}, \emph{robust}, \emph{model} --- luôn tìm định nghĩa tác giả đưa ra, hoặc tìm cách họ \emph{đo} khái niệm đó ở phần phương pháp.}
\ar[2]{Vì sao học từ rời ngữ cảnh kém hiệu quả}{Bạn nhớ tên mà không nhớ ``chữ ký'': không biết từ đó đi với giới từ nào, danh từ nào, ở văn vực nào.}
\ar[2]{Tam giác nghĩa}{Từ \(\rightarrow\) khái niệm trong đầu người dùng \(\rightarrow\) sự vật. Hiểu nhầm kỹ thuật thường nằm ở đỉnh khái niệm.}
\ar[2]{Cụm giới hạn phạm vi}{\emph{in this setting}, \emph{for small \(n\)}, \emph{under the assumption that} --- hay bị lướt qua và là nguyên nhân chính khiến người đọc hiểu kết quả mạnh hơn ý tác giả.}
\ar[1]{Vì sao dùng từ điển đơn ngữ cho người học}{Nó định nghĩa bằng ngôn ngữ đơn giản, kèm ví dụ thật, kèm giới từ và kết hợp từ --- những thứ từ điển song ngữ bỏ qua.}
\end{onbai}

\begin{hoidap}
\qa{Tôi nên bỏ hẳn từ điển Anh--Việt à?}{Không. Hãy dùng đúng vai trò: từ điển song ngữ tốt cho việc \emph{nhận ra} nghĩa khi đang đọc nhanh và không muốn đứt mạch. Từ điển đơn ngữ dành cho người học\cite{oald,ldoce} thì bắt buộc cho những từ bạn định \emph{dùng lại} --- vì bạn cần ví dụ, giới từ đi kèm và văn vực. Quy tắc thực dụng: tra song ngữ khi đọc, tra đơn ngữ khi ghi vào sổ từ.}
\qa{Làm sao biết một từ thuộc văn vực nào?}{Ba dấu hiệu. Một, nhãn trong từ điển học tập --- \emph{formal}, \emph{informal}, \emph{technical}, \emph{spoken}. Hai, nguồn gốc: từ gốc Latin hoặc Hy Lạp thường trang trọng hơn từ gốc Germanic cùng nghĩa (\emph{purchase} so với \emph{buy}, \emph{commence} so với \emph{start}); chi tiết ở \ptr{A/04}. Ba, và đáng tin nhất: tra trong kho ngữ liệu\cite{coca} rồi xem từ đó xuất hiện nhiều ở loại văn bản nào --- báo chí, học thuật, hay hội thoại.}
\qa{Trong bài báo, khi nào tôi nên viết \emph{show} và khi nào \emph{suggest}?}{Dùng \emph{show} khi dữ liệu trong chính bài của bạn đủ để kết luận trực tiếp --- một hình vẽ, một bảng số, một chứng minh. Dùng \emph{suggest} hoặc \emph{indicate} khi bằng chứng có tính gián tiếp, khi mẫu nhỏ, hoặc khi kết luận vượt ra ngoài điều kiện thí nghiệm. Viết mạnh hơn mức bằng chứng cho phép là lỗi nghiêm trọng hơn viết dè dặt: người phản biện đọc ra ngay, và nó làm giảm độ tin cậy của cả bài. Chi tiết ở \ptr{C/14}.}
\qa{Tôi hay hiểu nhầm những câu dài trong bài báo. Đó là vấn đề từ vựng hay ngữ pháp?}{Thường là ngữ pháp ở mức \emph{cấu trúc mệnh đề}, không phải từ vựng. Câu học thuật dài hay có ba tầng: mệnh đề chính, các mệnh đề phụ bổ nghĩa, và các cụm giới hạn phạm vi. Nếu bạn tách được mệnh đề chính ra trước thì phần còn lại trở nên dễ. Chương \ptr{A/03} dạy đúng kỹ năng tách đó, và nó là lượng ngữ pháp lớn nhất mà việc đọc thật sự đòi hỏi.}
\qa{Có cách nào nhanh để kiểm tra một cụm tiếng Anh có tự nhiên không?}{Tra trong kho ngữ liệu và so tần suất giữa hai cách viết bạn đang cân nhắc. Ví dụ, so \emph{make a decision} với \emph{do a decision}, hay so \emph{significant improvement} với \emph{big improvement} trong văn học thuật. Điều cần nhớ: tần suất nói lên tính tự nhiên, không nói lên tính đúng ngữ pháp --- và trong viết học thuật, tự nhiên mới là thứ bạn đang nhắm tới. Cách làm cụ thể ở \ptr{E/20}.}
\qa{Tôi nên ghi sổ từ vựng như thế nào cho đúng?}{Đừng ghi cặp từ--nghĩa. Mỗi mục nên có: câu thật bạn gặp từ đó (giữ nguyên ngữ cảnh), nghĩa bằng tiếng Anh đơn giản, hai hoặc ba kết hợp từ hay gặp, và nhãn văn vực nếu có. Chương \ptr{A/02} nói vì sao bốn thành phần này là tối thiểu, và \ptr{E/20} nói cách biến chúng thành thẻ ghi nhớ.}
\qa{Khi đọc, tôi nên tra từ ngay hay đọc tiếp rồi tra sau?}{Tuỳ mục tiêu của lần đọc đó. Nếu đang đọc để \emph{hiểu nội dung}, hãy đọc tiếp và chỉ tra những từ chặn hẳn mạch hiểu --- việc đoán nghĩa từ ngữ cảnh vừa nhanh hơn vừa là kỹ năng cần luyện (\ptr{A/04}). Nếu đang đọc để \emph{học ngôn ngữ}, hãy chọn một đoạn ngắn và tra kỹ mọi thứ, kể cả giới từ và kết hợp. Sai lầm phổ biến là trộn hai mục tiêu: tra từng từ trong một bài dài, mất mạch, và cuối cùng không nhớ nội dung lẫn từ mới.}
\end{hoidap}

\begin{baitap}
\bt[1]{Lấy một đoạn tóm tắt bài báo trong ngành bạn, gạch chân mọi từ chỉ mức cam kết}{Bài báo bất kỳ}{\emph{show, suggest, may, appear, tend to}\ldots{} rồi viết một câu nói tác giả cam kết mạnh tới đâu.}
\bt[1]{Tra ba từ bạn tưởng đã biết trong từ điển học tập\cite{oald}}{Từ điển đơn ngữ}{Ghi lại những gì bạn \emph{không} biết trước đó: giới từ đi kèm, tính đếm được, văn vực.}
\bt[2]{So sánh \emph{cause}, \emph{lead to}, \emph{result in} trong kho ngữ liệu}{\cite{coca}}{Ghi lại năm danh từ hay đi sau mỗi từ; rút ra kết luận về sắc thái.}
\bt[2]{Áp bốn phép kiểm ở mục 4 cho ba câu khó trong tài liệu bạn đang đọc}{Tài liệu của bạn}{Đặc biệt làm kỹ phép xác định phạm vi --- ghi ra câu đó \emph{không} áp dụng cho trường hợp nào.}
\bt[2]{Viết lại ví dụ ở đầu chương thành hai phiên bản: một cam kết mạnh, một dè dặt hơn}{Bài tập viết}{So hai phiên bản và giải thích mỗi thay đổi từ ngữ làm đổi điều gì.}
\bt[3]{Tìm ba từ trong ngành của bạn có nghĩa khác với nghĩa đời thường}{Tài liệu chuyên ngành}{Ví dụ: \emph{significant}, \emph{robust}, \emph{noise}. Viết định nghĩa thao tác cho mỗi từ như tác giả sẽ viết trong bài báo.}
\bt[3]{Chọn một đoạn văn bạn đã viết bằng tiếng Anh và soi lại theo ba tầng nghĩa}{Bài viết của bạn}{Với mỗi động từ báo cáo kết quả, hỏi: mức cam kết này có đúng với bằng chứng tôi có không?}
\end{baitap}

\section*{Kết chương và đọc tiếp}

Chương này thay giả định ``mỗi từ một nghĩa'' bằng một mô hình ba tầng, và đưa ra bốn phép kiểm để biết mình đã hiểu đúng thay vì chỉ thấy quen. Hai điều đáng mang theo: thang cam kết của các động từ báo cáo, vì nó là chỗ chứa nhiều thông tin nhất trong văn khoa học; và thói quen tìm cụm giới hạn phạm vi trước khi kết luận tác giả nói gì.

Nhưng nếu nghĩa nằm ở cách dùng, thì ``học một từ'' nghĩa là gì, và cần bao nhiêu từ để đọc được một bài báo mà không phải tra liên tục? Chương \ptr{A/02} trả lời cả hai câu bằng những con số cụ thể --- và các con số đó sẽ định hình toàn bộ lộ trình luyện tập ở phần E.
\end{document}
